\documentclass[aspectratio=169]{beamer}
\usetheme{Madrid}
\usecolortheme{dolphin}
\setbeamertemplate{navigation symbols}{}
\setbeamertemplate{footline}[frame number]
\setbeamersize{text margin left=6mm,text margin right=6mm}
\usepackage[T1]{fontenc}
\usepackage{lmodern}
\usepackage{booktabs,tabularx,array}
\usepackage{xcolor}
\IfFileExists{xurl.sty}{\usepackage{xurl}}{\usepackage{url}}
\hypersetup{hidelinks,pdftitle={Attempt 4 - Final Query Result}}
\newcommand{\good}[1]{\textcolor{green!45!black}{\textbf{#1}}}
\newcommand{\bad}[1]{\textcolor{red!75!black}{\textbf{#1}}}
\title{Attempt 4: Final Query Result}
\subtitle{The two-dimensional model works locally, but not consistently enough}
\author{}
\date{2026-07-29}

\begin{document}
\begin{frame}\titlepage\end{frame}

\begin{frame}{1. The Answer First}
\begin{alertblock}{Final decision}
The current shared two-dimensional model is \textbf{not} a general
replacement for product quantization.
\end{alertblock}
\begin{itemize}
  \item It has clear wins in some settings.
  \item It loses or ties in other equally important settings.
  \item No storage budget passes both datasets and both index settings.
\end{itemize}
\begin{block}{What this means}
The implementation works and the idea has a real signal, but the signal is
not robust enough for a SIGMOD/VLDB/ICDE contribution in its current form.
\end{block}
\end{frame}

\begin{frame}{2. What Was the Question?}
We compress every database vector to either \textbf{32 bytes} or
\textbf{64 bytes}.
\medskip
\begin{columns}[T]
\column{0.47\textwidth}
\begin{block}{Standard approach}
Split coordinates into small blocks. Learn a separate set of representative
points for every block.
\end{block}
\column{0.47\textwidth}
\begin{block}{Our approach}
Learn one reusable two-dimensional shape. Each coordinate pair stretches,
rotates, and shifts that shared shape.
\end{block}
\end{columns}
\begin{alertblock}{The test}
At the same code bytes and the same searched database candidates, does our
model recover more true neighbours without slowing queries?
\end{alertblock}
\end{frame}

\begin{frame}[fragile]{3. How the Method Works}
\small
\begin{verbatim}
Training:
  learn one shared 2D shape
  for each coordinate pair:
    learn a shift + 2-by-2 transform
  store one label per database pair
Query:
  build 64 small tables per searched cell
  scan labels and keep the best 100 IDs
\end{verbatim}
\textbf{Why it might help:} a non-rectangular shape can follow natural
two-dimensional geometry better than two independent scalar grids.
\medskip

\textbf{Why it might fail:} every searched cell needs new tables, and one
shared shape may be too restrictive for all datasets and rates.
\end{frame}

\begin{frame}{3A. Running Example: Reuse One Shape}
\small
For arithmetic only, use four shared points instead of the real 16 or 256:
\[
z_0=(-1,0),\quad z_1=(0,-1),\quad
z_2=(1,0),\quad z_3=(0,1).
\]
Each coordinate pair applies its own
\(\text{centre}_{g,k}=\mu_g+A_gz_k\):
\begin{center}
\begin{tabular}{c c c c}
\toprule
label \(k\) & shared \(z_k\) & pair A: stretch + shift & pair B: rotate + shift\\
\midrule
0 & \((-1,0)\) & \((-1,1)\) & \((-1,-1)\)\\
1 & \((0,-1)\) & \((1,0)\)  & \((0,0)\)\\
2 & \((1,0)\)  & \((3,1)\)  & \((-1,1)\)\\
3 & \((0,1)\)  & \((1,2)\)  & \((-2,0)\)\\
\bottomrule
\end{tabular}
\end{center}
\begin{block}{What is shared?}
Both pairs reuse the four labels and the same underlying shape. Only a
two-number shift and a \(2\times2\) transform change by pair.
\end{block}
\end{frame}

\begin{frame}{3B. Running Example: Build One Query Table}
\small
For pair A, let the transformed query residual be \(q=(2.8,1.2)\).
Build four squared-distance entries:
\begin{center}
\begin{tabular}{c c c}
\toprule
label & reconstructed centre & \(T[k]=\lVert q-\text{centre}_k\rVert^2\)\\
\midrule
0 & \((-1,1)\) & \(14.48\)\\
1 & \((1,0)\)  & \(4.68\)\\
2 & \((3,1)\)  & \(\mathbf{0.08}\)\\
3 & \((1,2)\)  & \(3.88\)\\
\bottomrule
\end{tabular}
\end{center}
A database vector storing label \(2\) for this pair contributes
\(\mathbf{T[2]=0.08}\); its coordinates are not decoded during the scan.
\begin{alertblock}{Scale up to the real method}
A full code stores 64 labels. Its score is
\(T_0[k_0]+\cdots+T_{63}[k_{63}]\): 64 table lookups. Labels use 4 or 8 bits,
so the vector payload is exactly 32 or 64 bytes.
\end{alertblock}
\end{frame}

\begin{frame}{4. Fair Comparison}
\small
\begin{tabularx}{\textwidth}{>{\bfseries}p{0.29\textwidth} X}
data & standard SIFT1M and GIST1M queries and exact answers\\
quality & Recall@100: fraction of the true 100 neighbours recovered\\
speed & complete queries per second and single-query tail latency\\
storage & exactly 32 or 64 code bytes per database vector\\
coarse index & 1,024 and 4,096 database cells\\
search range & nine fixed numbers of cells searched per query\\
repetition & one warmup and seven measured runs\\
hardware & one CPU socket; one-thread latency and 12-core throughput\\
\end{tabularx}
\normalsize
\begin{block}{Included in time}
Query transformation, cell selection, table construction, database scan, and
selection of the best 100 IDs.
\end{block}
\end{frame}

\begin{frame}{5. What Counts as a Real Win?}
For one operating point, either:
\begin{enumerate}
  \item \textbf{speed route}: at least 10\% more queries per second at nearly
  the same Recall, without making p95 latency more than 5\% worse; or
  \item \textbf{quality route}: at least 0.002 more Recall at nearly the same
  throughput.
\end{enumerate}
\medskip
It must also:
\begin{itemize}
  \item use no more complete index bytes;
  \item use no more than twice the baseline construction CPU; and
  \item not be beaten simultaneously in quality, speed, and bytes by the
  contextual comparison.
\end{itemize}
\begin{alertblock}{Terminal rule}
One byte budget and one route must pass \textbf{both datasets and both coarse
index sizes}. A single attractive point is not enough.
\end{alertblock}
\end{frame}

\begin{frame}{6. The Complete Matrix Is Valid}
\begin{center}
\begin{tabular}{lrr}
\toprule
Check & Completed & Required\\
\midrule
logical method settings & 94 & 94\\
warmup passes & 188 & 188\\
measured passes & 1,316 & 1,316\\
CPU hours & 239.25 & at most 256\\
wall hours & 101.22 & at most 120\\
peak memory & 11.12 GiB & at most 16 GiB\\
\bottomrule
\end{tabular}
\end{center}
\begin{block}{Stability}
For every fixed point, all repetitions returned the same IDs, candidate count,
and Recall. A second summary run reproduced every output byte-for-byte.
\end{block}
\scriptsize One interrupted warmup remained conservatively charged and was
later rerun successfully; it is not a scientific failure.
\end{frame}

\begin{frame}{7. Final Result by Dataset and Index Size}
\small
\begin{center}
\begin{tabular}{llcc}
\toprule
Dataset & Coarse cells & 32 bytes & 64 bytes\\
\midrule
GIST & 1,024 & \good{pass} & \good{pass}\\
GIST & 4,096 & \bad{fail: index bytes} & \bad{fail: index bytes}\\
SIFT & 1,024 & \bad{fail} & \bad{fail}\\
SIFT & 4,096 & \bad{fail} & \good{pass}\\
\bottomrule
\end{tabular}
\end{center}
\normalsize
\begin{alertblock}{Why the project-level answer is no}
Neither column is all green. The 32-byte model loses on SIFT. The 64-byte
model helps SIFT only with 4,096 coarse cells, not with 1,024.
\end{alertblock}
\end{frame}

\begin{frame}{8. The Clearest Positive Point}
\begin{block}{SIFT, 4,096 coarse cells, 64-byte code}
Searching 256 cells:
\[
\text{Recall}=0.9261,\qquad \text{throughput}=966\text{ queries/s}.
\]
\end{block}
Compared with the eligible full-dimensional OPQ point:
\begin{itemize}
  \item about \good{56\% more throughput} at matched Recall;
  \item about \good{29\% lower p95 latency}; and
  \item quality advantage of \good{0.00214} at matched throughput.
\end{itemize}
\begin{alertblock}{Why this still does not establish the contribution}
The benefit appears only at the 64-byte, 4,096-cell, high-search setting.
The same model does not repeat it at 1,024 cells or at 32 bytes.
\end{alertblock}
\end{frame}

\begin{frame}{9. The Decisive Negative Evidence Is SIFT}
\begin{itemize}
  \item At 32 bytes and 1,024 cells, S's best raw speed comparison is about
  \bad{28\% slower}; a representative quality comparison is about
  \bad{0.079 Recall worse}.
  \item At 32 bytes and 4,096 cells, its best raw speed comparison is about
  \bad{19\% slower}, and the best quality comparison is about
  \bad{0.042 Recall worse}.
  \item At 64 bytes and 1,024 cells, quality is nearly tied but slightly
  negative; S is about \bad{11\% slower} with about
  \bad{17\% worse p95 latency}.
\end{itemize}
\begin{block}{Interpretation}
The shared shape is useful only after enough code capacity and enough scanned
candidates amortize its tables. That is an operating-region effect, not a
robust new frontier.
\end{block}
\end{frame}

\begin{frame}{10. Why GIST Looks Better---and Why We Are Careful}
GIST has 960 dimensions, but S stores a label only for the first 128
transformed coordinates.
\medskip
For the remaining 832 coordinates, every vector in one coarse cell reuses the
cell centre.
\begin{center}
\textbf{more searched cells}
$\rightarrow$ more approximate tail values
$\rightarrow$ possible false positives
\end{center}
\begin{block}{What is valid}
The GIST numbers are valid for the system definition fixed before seeing
queries.
\end{block}
\begin{alertblock}{What is not valid}
They are not clean proof that the shared 2D shape alone is better. At 4,096
cells, the extra routing state also makes S's complete index larger than the
matched baseline.
\end{alertblock}
\end{frame}

\begin{frame}{11. Scientific Result vs. Engineering Result}
\begin{columns}[T]
\column{0.47\textwidth}
\begin{block}{Scientific evidence}
\begin{itemize}
  \item Shared 2D geometry can improve real query ordering.
  \item The gain depends strongly on rate and candidate regime.
  \item Offline reconstruction gains do not predict a common query win.
\end{itemize}
\end{block}
\column{0.47\textwidth}
\begin{block}{Artifact engineering}
\begin{itemize}
  \item deterministic packed scanner;
  \item shared candidate schedule;
  \item complete byte and resource accounting;
  \item 101-hour reproducible matrix.
\end{itemize}
\end{block}
\end{columns}
\begin{alertblock}{Reviewer view}
The engineering makes the negative conclusion trustworthy. It does not turn
the local positive points into a general research contribution.
\end{alertblock}
\end{frame}

\begin{frame}{12. Final Takeaway and Next Question}
\begin{alertblock}{Current direction}
\textbf{Close it under the frozen hypothesis.} Do not rescue it by changing
datasets, probe grids, thresholds, or storage accounting after seeing results.
\end{alertblock}
\begin{block}{Evidence worth preserving}
The 64-byte high-candidate SIFT region is real. It shows that the mechanism is
not empty; it is simply not robust.
\end{block}
\begin{block}{A scientifically different next question}
Why does the shared shape help only when code capacity and candidate count are
high, and can that mechanism be made robust without a query-selected policy
or extra per-cell state?
\end{block}
\end{frame}

\begin{frame}[allowframebreaks]{References and Evidence}
\tiny
\begin{itemize}
  \item Jégou, Douze, and Schmid. ``Product Quantization for Nearest Neighbor
  Search.'' TPAMI 2011.
  \item Ge et al. ``Optimized Product Quantization.'' CVPR 2013.
  \item André, Kermarrec, and Le Scouarnec. ``Quicker ADC.'' TPAMI.
  \item Gao and Long. ``RaBitQ.'' SIGMOD 2024.
  \item Gao et al. ``Extended RaBitQ.'' 2024.
  \item Li et al. ``SAQ.'' 2025.
  \item Frozen design:
  \texttt{structured\_2d\_fair\_query\_evaluation\_2026\_07\_24.md}.
  \item Final result:
  \texttt{structured\_2d\_natural\_query\_result\_2026\_07\_29.md}.
\end{itemize}
\end{frame}

\end{document}
